\documentclass[11pt]{article}
\usepackage[margin=1.1in]{geometry}
\usepackage{amsmath, amssymb, amsthm}
\usepackage{booktabs}
\usepackage{xcolor}
\definecolor{linknavy}{RGB}{25, 55, 125}
\usepackage[colorlinks=true, linkcolor=linknavy, citecolor=linknavy, urlcolor=linknavy]{hyperref}
\usepackage{orcidlink}
\usepackage{fancyhdr}

\newtheorem{theorem}{Theorem}[section]
\newtheorem{proposition}[theorem]{Proposition}
\newtheorem{lemma}[theorem]{Lemma}
\newtheorem{conjecture}[theorem]{Conjecture}
\theoremstyle{definition}
\newtheorem{definition}[theorem]{Definition}
\theoremstyle{remark}
\newtheorem{remark}[theorem]{Remark}

\newcommand{\Z}{\mathbb{Z}}
\newcommand{\R}{\mathbb{R}}
\newcommand{\C}{\mathbb{C}}
\newcommand{\Qp}{\mathbb{Q}_p}
\newcommand{\zmax}{z_{\max}}

\title{The Bottom of the Shub--Smale Tau Conjecture:\\
an Exact Census of Integer Roots for Constant-Free\\
Straight-Line Programs of Length at Most Seven}
\author{Felipe Santiba\~nez-Leal\,\orcidlink{0000-0002-0150-3246}\\
\small CAOS Research program, \texttt{github.com/fsantibanezleal/CAOS\_RESEARCH}}
\date{2026-08-02\\[2pt]\normalsize Version 0.02}

\pagestyle{fancy}
\fancyhf{}
\fancyhead[L]{\small CAOS Research Preprint}
\fancyhead[R]{\small Tau-Conjecture Census \textperiodcentered\ v0.02}
\fancyfoot[C]{\small \thepage}
\renewcommand{\headrulewidth}{0.4pt}
\fancypagestyle{plain}{\fancyhf{}\fancyfoot[C]{\small \thepage}\renewcommand{\headrulewidth}{0pt}}

\begin{document}
\maketitle
\begin{center}\small
\fbox{\parbox{0.92\textwidth}{\centering
\textbf{PREPRINT} \textperiodcentered\ not peer reviewed \textperiodcentered\ CAOS Research programme\\
First paper of the tau-conjecture record \textperiodcentered\ version 0.02 \textperiodcentered\ 2026-08-02\\
DOI \href{https://doi.org/10.5281/zenodo.21753438}{10.5281/zenodo.21753438} (concept DOI, resolves to the latest version)\\
Licensed CC BY 4.0 \textperiodcentered\ sources and machine records: \texttt{github.com/fsantibanezleal/CAOS\_RESEARCH}
}}
\end{center}

\begin{abstract}
For a univariate integer polynomial $f$, let $\tau(f)$ be the minimum number of
$+,-,\times$ gates needed to compute $f$ from $x$ and the constants $-1,0,1$, and let
$z(f)$ be its number of distinct integer roots. The Shub--Smale tau conjecture
(Smale's fourth problem) asserts $z(f) \le (1+\tau(f))^{\kappa}$ for a universal
$\kappa$; it implies $\mathrm{P}_{\C} \ne \mathrm{NP}_{\C}$ in the Blum--Shub--Smale
model and $\mathrm{VP}^0 \ne \mathrm{VNP}^0$, and it is open even for $\kappa = 1$.
We report the first exact census of the conjecture's growth function at the bottom of
its ladder: writing $\zmax(\tau)$ for the maximum of $z(f)$ over all nonzero $f$ with
$\tau(f) \le \tau$, we prove by exhaustive, exactly verified computation that
\[
\zmax(1),\dots,\zmax(7) \;=\; 1,\, 2,\, 3,\, 3,\, 4,\, 5,\, 5 .
\]
In particular the minimum cost of $4$ distinct integer roots is $5$ gates, of $5$
roots is $6$ gates, and (new in this version) of $6$ roots is exactly $8$ gates, by an
explicit verified witness: $q(q-2)(q-6)$ for $q = x^2 - x$, computable in $8$ gates by
chained subtraction sharing; the growth function has plateaus at $\tau = 4$ and
$\tau = 7$, so an extra gate does not always buy an extra root, and
$z_{\max}(8) \ge 6$. The censuses are decision-complete: depth $6$ required the exact construction of
all $25{,}844{,}905$ reachable computation states, and depth $7$ a complete scan of
$2{,}013{,}706$ new polynomials, using a ``last-gate'' lemma that decides one depth
beyond any exhausted frontier without storing the next frontier. The enumerator is
anchored on Markstr\"om's published exhaustive census of integer targets (all
fourteen anchor values reproduced exactly) and cross-checked against an independent
computer-algebra root counter on $284$ polynomials. The six-root threshold was
closed by a case split: the last gate of an $8$-gate $6$-rooter must involve the
seventh value, and if it is a multiplication its root count is a union of two known
root sets, so that case reduces to a co-occurrence query over the depth-$6$ frontier;
the scan found $408$ witnesses. We further prove two elementary
stall theorems explaining the observed record mechanisms: for any monic
$h \in \Z[x]$ of degree at least $2$, towers built by iterating $h$ alone have
depth-independent integer root counts (for $h = x^2 - 2$, the iterates $h^{\circ k}(x) - x$
keep exactly $2$ integer roots against $2^k$ real roots), and across the quadratic
family $h_c = x^2 - c$ with $c \le 200$ the maximum tower yield is $5$, attained only
at $c = 2$; a second yield series at $c = m^2 + m + 1$, produced by genuine integer
$2$-cycles, is identified and closed by the classical fact that integer polynomial
cycles have length at most $2$. All computations are exact integer arithmetic,
reproducible from committed code with hypotheses declared before each run; three of our
own pre-registered predictions were refuted by the machine and are reported as such.
This is an experimental-mathematics record in the tradition of Markstr\"om's census:
it decides nothing about the conjecture asymptotically, and says so.
\end{abstract}

\section{Introduction}

\subsection{The conjecture}

A \emph{constant-free straight-line program} (SLP) over $\Z[x]$ is a sequence
$(1, x, f_2, \dots, f_N)$ in which every $f_i$ ($i \ge 2$) is a sum, difference, or
product of two earlier entries; following \cite{ShubSmale1995, Rojas2003} we write
$\tau(f)$ for the least $N - 1$ over all such programs with $f_N = f$. Equivalently
\cite{Burgisser2024}, $\tau(f)$ is the minimal number of operation gates of an
arithmetic circuit over the inputs $x$ and the free constants $-1, 0, 1$. Let $z(f)$
denote the number of distinct roots of $f$ in $\Z$.

\begin{conjecture}[Shub--Smale \cite{ShubSmale1995}; Smale's fourth problem \cite{Smale1998}]
There is a universal constant $\kappa \ge 1$ such that every nonzero
$f \in \Z[x]$ satisfies $z(f) \le (1 + \tau(f))^{\kappa}$.
\end{conjecture}

The conjecture fails for $\kappa < 1$: the geometric-progression products
$(x - 2)(x - 4)\cdots(x - 2^{2^j})$ acquire one root per $O(1)$ gates
\cite{Rojas2003}. It is open even for $\kappa = 1$. Its significance is structural:
Shub and Smale proved that it implies
$\mathrm{P}_{\C} \ne \mathrm{NP}_{\C}$ in the Blum--Shub--Smale model
\cite{ShubSmale1995, BCSS1998}, the proof reducing to the hardness of computing all
sequences $(m_n \cdot n!)$; B\"urgisser \cite{Burgisser2009}, building on Koiran
\cite{Koiran2004}, proved that it implies $\mathrm{VP}^0 \ne \mathrm{VNP}^0$, and
that the single sequence $(n!)$ being hard to compute already yields that separation.
Easy factorials would give nonuniform polynomial-time integer factoring
\cite{BCSS1998, Cheng2003}; with integer division allowed, $n!$ \emph{is} easy
\cite{Shamir1979}, so the restriction to $+,-,\times$ carries the entire content. The
real-zeros analogue of the conjecture is false: iterating the logistic map produces
$2^j$ real roots at cost $O(j)$ \cite{Rojas2003}, which motivated Koiran's real tau
conjecture for structured expressions \cite{Koiran2011, KPTT2015, Hrubes2013,
BriquelBurgisser2020, Dutta2021}. On the lower-bound side, no nontrivial lower bound
on $\tau(n!)$ is known \cite{Koiran2004}; on the reduction side, Rojas proved that a
$p$-adic digit version of the conjecture, for any fixed prime, already implies the
full statement, via a $p$-independent subquadratic bound on the valuation spectrum of
roots \cite{Rojas2003}. Recent work applies the conjecture to explicit exponential
lower bounds \cite{BBDM2026} and tightens the bridge from BSS separations to
constant-free Valiant separations \cite{Burgisser2026}; the current survey reference
is \cite[\S4.6]{Burgisser2024}.

\subsection{What this paper does}

The literature contains, to our knowledge, exactly one exhaustive computation in this
model: Markstr\"om's census of \emph{integer targets} (the straight-line complexity
of small factorials and primorials), complete to program length $11$ on 2013-era
hardware \cite{Markstrom2014}. No published computation addresses the
\emph{polynomial} side, i.e.\ the actual growth function of the conjecture. We define
\[
\zmax(\tau) \;=\; \max\{\, z(f) \;:\; f \in \Z[x] \setminus \{0\},\ \tau(f) \le \tau \,\}
\]
and determine it exactly for $\tau \le 7$, together with the full record galleries,
explicit optimal witnesses, and the exact cost thresholds for $3, 4, 5$ and (within
one gate) $6$ distinct integer roots. We also prove two elementary theorems that
explain the observed record mechanisms and close entire construction classes
(Section~\ref{sec:stall}), and we measure the quadratic tower family exhaustively
(Section~\ref{sec:family}).

Our contributions:
\begin{enumerate}
\item \textbf{The census} (Section~\ref{sec:census}): $\zmax(1..7) = 1,2,3,3,4,5,5$,
decision-complete, with per-depth state and polynomial counts, record galleries and
replayed witness programs. Plateaus occur at $\tau = 4$ and $\tau = 7$.
\item \textbf{Exact thresholds}: the minimum $\tau$ admitting $3$ distinct integer
roots is $3$ ($x^3 - x$); for $4$ roots it is $5$; for $5$ roots it is $6$; for $6$
roots it is exactly $8$ (Section~\ref{sec:window}: the census excludes $7$, and the
scan of this version produced explicit verified $8$-gate witnesses).
\item \textbf{A method lemma} (Lemma~\ref{lem:lastgate}): every polynomial of cost
$d + 1$ is one gate over a normalized depth-$d$ state, so $\zmax(d+1)$ is computable
from an exhausted depth-$d$ frontier without storing the depth-$(d+1)$ frontier. This
bought depth $7$ (a $\sim$$10^9$-state frontier avoided).
\item \textbf{Stall theorems} (Section~\ref{sec:stall}): iterating any fixed monic
map of degree $\ge 2$ yields depth-independent integer root counts; for the
Chebyshev-conjugate map $x^2 - 2$, whose difference-of-squares splittings generate
all census records at depths $5$--$7$, the exact stable core is $\{0, \pm 1, \pm 2\}$.
\item \textbf{The family measurement} (Section~\ref{sec:family}): across
$h_c = x^2 - c$, $c \le 200$, and all tower shapes to four iterations, the maximum
integer-root yield is $5$, attained only at $c = 2$; two yield-$4$ series exist,
$c = m(m+1)$ (fixed and anti-fixed points) and $c = m^2 + m + 1$ (genuine
$2$-cycles), and the classical bound on integer polynomial cycle lengths shows this
inventory is complete.
\item \textbf{Anchoring and audit} (Section~\ref{sec:anchor}): the enumerator
reproduces all fourteen of Markstr\"om's published anchor values; an independent
computer-algebra cross-check agrees on $284$ polynomials; every experiment's
hypothesis, budget and kill criterion were committed before its run, and two of our
pre-registered predictions were refuted by the machine (both disclosed below).
\end{enumerate}

Every claim in this paper is labeled: \textbf{[MV]} machine-verified in exact
arithmetic from committed code, \textbf{[D]} proved here (elementary, self-contained),
or \textbf{[C]} conjectural/observational. Nothing here decides the conjecture in
either direction, and no asymptotic consequence is claimed: the census maps where the
conjecture lives at the bottom of its ladder, in the same genre as
\cite{Markstrom2014}.

\section{Model, normalization, and the census method}\label{sec:model}

\begin{definition}\label{def:model}
Programs operate on values in $\Z[x]$. The inputs are $-1$, $1$ and $x$; each gate
applies $+$, $-$ or $\times$ to two (not necessarily distinct) earlier values; the
length of a program is its number of gates; $\tau(f)$ is the minimal length of a
program computing $f$ at some gate.
\end{definition}

\begin{lemma}[free $0$ is redundant]\label{lem:zero}
Adding $0$ as a fourth free input does not change $\tau$. \emph{[D]}
\end{lemma}
\begin{proof}
In any program using the input $0$: $a + 0$ and $a - 0$ duplicate $a$; $a \times 0$
computes $0$; $0 - a$ equals $(-1) \times a$, one gate either way. Replacing each use
and deleting duplicated gates never lengthens the program.
\end{proof}

\begin{lemma}[normalization]\label{lem:norm}
Every $f$ has an optimal program in which no gate recomputes an input or an earlier
value, and no gate computes $0$. \emph{[D]}
\end{lemma}
\begin{proof}
A gate duplicating an available value can be deleted, rewiring its uses to the
original; a computed $0$ is usable only in the redundant patterns of
Lemma~\ref{lem:zero}. Each deletion shortens the program, so an optimal one contains
neither.
\end{proof}

\begin{lemma}[reached-set sufficiency]\label{lem:state}
The extensions available to a program depend only on the \emph{set} of values it has
computed. Consequently a breadth-first search over reached sets (``states''), with
set-level deduplication, visits every normalized reachable state of each depth
exactly once. \emph{[D]}
\end{lemma}

This is the polynomial analogue of the range-isomorphism reduction of
\cite{Markstrom2014}. The census algorithm is exactly this BFS, in exact integer
arithmetic on dense coefficient tuples, with each new polynomial's first appearance
depth recorded (that depth is $\tau(f)$, by Lemmas~\ref{lem:norm} and
\ref{lem:state}). Integer roots are counted exactly: strip the power of $x$, then
test the divisors of the trailing coefficient (both signs).

\begin{lemma}[last gate]\label{lem:lastgate}
Let $\mathcal{F}_d$ be the set of normalized reached states of depth $d$. Every $f$
with $\tau(f) = d + 1$ is $a \circ b$ for a single gate $\circ$ and operands $a, b$
available in some state of $\mathcal{F}_d$. Hence
$\zmax(d+1) = \max\bigl(\zmax(d),\ \max z(a \circ b)\bigr)$ over one op over every
state of $\mathcal{F}_d$, and this maximum is computable without constructing
$\mathcal{F}_{d+1}$. \emph{[D]}
\end{lemma}
\begin{proof}
Take an optimal program of length $d + 1$ for $f$ and normalize it
(Lemma~\ref{lem:norm}; deletions would shorten it, contradicting optimality, so it is
already normalized). Its first $d$ gates form a normalized reached state of depth
$d$, and its last gate is one op over that state's available values.
\end{proof}

\section{Anchoring and audit discipline}\label{sec:anchor}

\textbf{Markstr\"om anchor [MV].} Restricted to the integer model (start value $1$,
positive normalized values), our enumerator reproduces Figure~1 of
\cite{Markstrom2014} exactly at every compared depth: reached-set sizes
$2, 4, 9, 26, 102, 562, 4363$ and complete initial intervals
$2, 4, 6, 12, 40, 112, 310$ for lengths $1$ through $7$: fourteen of fourteen anchor
values.

\textbf{Independent cross-check [MV].} An independent root counter (SymPy's exact
factorization) agrees with ours on all $284$ polynomials checked: every polynomial of
$\tau \le 3$ and every record polynomial of the censuses.

\textbf{Engine re-derivation [MV].} The depth-$7$ computation uses a reimplemented,
memoized engine; before its new output was trusted, it was required to re-derive every
prior census value (state counts $9 / 98 / 1462 / 29506 / 778087$, new-polynomial
counts through depth $6$, and the record sets), and did.

\textbf{Pre-registration.} Every computation ran under a hypothesis file committed to
the repository \emph{before} the run, with falsifiable predictions, an explicit
statement of what a pass and a fail would prove, a compute budget and a kill
criterion. Three committed predictions were refuted by the machine: we predicted
$\zmax(6) = 4$ (the census returned $5$; Section~\ref{sec:census}); we predicted zero
tower yield off the fixed-point series (a second series exists;
Section~\ref{sec:family}); and we predicted the six-root multiplication case empty at
$8$ gates (the scan found $408$ witnesses; Section~\ref{sec:window}). Each refutation
is reported alongside what it taught.
One tooling incident is likewise on record: divisor-based root counting is infeasible
against tower constant terms of size $c^{2^k}$, and was replaced, for the family
measurement only, by direct evaluation over a window proved correct in advance
(Section~\ref{sec:family}), cross-checked against the divisor method where both run.

\section{The census}\label{sec:census}

\begin{theorem}[the census; MV]\label{thm:census}
In the model of Definition~\ref{def:model},
\[
\begin{array}{c|ccccccc}
\tau & 1 & 2 & 3 & 4 & 5 & 6 & 7\\\hline
\zmax(\tau) & 1 & 2 & 3 & 3 & 4 & 5 & 5
\end{array}
\]
and the censuses behind each value are decision-complete.
\end{theorem}

The supporting exhaustive counts are:

\begin{center}
\begin{tabular}{lrrrrrrr}
\toprule
$\tau$ & 1 & 2 & 3 & 4 & 5 & 6 & 7\\
\midrule
reached states & 9 & 98 & 1{,}462 & 29{,}506 & 778{,}087 & 25{,}844{,}905 & (not stored)\\
new polynomials & 9 & 34 & 177 & 1{,}249 & 11{,}377 & 134{,}494 & 2{,}013{,}706\\
$\zmax$ & 1 & 2 & 3 & 3 & 4 & 5 & 5\\
\bottomrule
\end{tabular}
\end{center}

Depth $7$ was decided by Lemma~\ref{lem:lastgate} over the exhausted depth-$6$
frontier; among its $2{,}013{,}706$ new polynomials the distinct-integer-root
histogram is $742{,}635$ rootless, $904{,}421$ with one, $299{,}488$ with two,
$63{,}903$ with three, $3{,}196$ with four, $63$ with five, and none with six or
more.

\begin{theorem}[cost thresholds; MV]\label{thm:thresholds}
The minimum $\tau$ admitting $k$ distinct integer roots is: $k=1{:}\,1$;
$k=2{:}\,2$; $k=3{:}\,3$; $k=4{:}\,5$; $k=5{:}\,6$; $k=6{:}\,8$.
\end{theorem}

The $k = 6$ lower bound is Theorem~\ref{thm:census} ($\zmax(7) = 5$); the upper
bound is the explicit $8$-gate witness of Section~\ref{sec:window}. Representative optimal witnesses, replayed gate by gate in
exact arithmetic from the committed artifacts:

\begin{center}
\small
\begin{tabular}{lll}
\toprule
$\tau$ & record & witness program\\
\midrule
3 & $x^3 - x$, roots $\{0,\pm1\}$ & $x{\cdot}x;\ x^2{-}1;\ x{\cdot}(x^2{-}1)$\\
4 & $-x(x{+}1)(x{+}2)$, roots $\{0,-1,-2\}$ & $-1{-}x;\ 1{-}({-}1{-}x);\ x({-}1{-}x);\ (x^2{+}x)\dots$\\
5 & $x^2 - (x^2{-}2)^2 = -(x^2{-}1)(x^2{-}4)$ & $-2;\ x^2;\ x^2{-}2;\ (x^2{-}2)^2;\ x^2 - (x^2{-}2)^2$\\
6 & $\mp x(x^2{-}1)(x^2{-}4)$, roots $\{0,\pm1,\pm2\}$ & the $\tau{=}5$ record, then $\times\, x$\\
\bottomrule
\end{tabular}
\end{center}

\textbf{Plateaus.} $\zmax$ pauses at $\tau = 4$ and $\tau = 7$. The record galleries
explain the asymmetry: the step $5 \to 6$ succeeds because multiplying by the input
$x$ costs one gate and adjoins the root $0$ to a record avoiding it; the steps
$3 \to 4$ and $6 \to 7$ fail because every candidate next root beyond the ``free''
set $\{0, \pm 1, \pm 2\}$ requires a built constant, and constants cost gates. The
$\tau = 5$ records (ten polynomials) and $\tau = 6$ records (four) are all
difference-of-squares splittings over the map $x^2 - 2$ (Section~\ref{sec:stall});
at $\tau = 7$ the sixty-three five-rooters appear, but no six-rooter.

\subsection{Closing the six-root window}\label{sec:window}

Version 0.01 of this paper left the six-root threshold in $[8, 9]$. It is $8$:

\begin{theorem}[MV, both directions]\label{thm:window}
There is an $8$-gate constant-free SLP whose output has $6$ distinct integer roots,
and none with $7$ gates. An explicit witness (replayed gate by gate in exact
arithmetic) is, with $q = x(x-1)$:
\[
x{-}1;\quad 2;\quad q = x\cdot(x{-}1);\quad 4;\quad q{-}2;\quad (q{-}2){-}4 = q{-}6;\quad
q(q{-}2);\quad f = q(q{-}2)\cdot(q{-}6),
\]
whose root set is $\{-2, -1, 0, 1, 2, 3\}$. Hence $\zmax(8) \ge 6$.
\end{theorem}

The method is a case split. By Lemma~\ref{lem:lastgate} at depth $8$, an $8$-gate
$6$-rooter is one gate over a depth-$7$ state $S \cup \{v\}$; since
$\zmax(7) = 5$, the final gate must involve the seventh value $v$. If that gate is a
multiplication, $f = v \cdot b$ and $z(f) = |R_v \cup R_b|$: a pure co-occurrence
condition on known root sets, decided exactly by one scan of the $25{,}844{,}905$
depth-$6$ states ($2$h$23$m single-core): $408$ hits, three reconstructed as explicit
programs and verified independently of the scan machinery. Two committed expectations
fell in the process: we had predicted this case empty (our schema hunt costed the
witness family at $9$ gates because it built the constants $\{2, 6\}$ independently;
the scan's witness builds $\{2, 4\}$ and reuses the subtraction chain
$q \mapsto q-2 \mapsto q-6$: a one-gate saving our cost model missed), and we had
predicted that all $\tau \le 7$ five-rooters have root set $\{0, \pm1, \pm2\}$:
in fact the $67$ five-rooters realize seven root-set patterns, including the
non-consecutive $\{-1,0,1,2,4\}$ and $\{-4,-2,-1,0,1\}$, and it is precisely such
shifted and punctured sets that leave room for a co-occurring factor to complete six
roots.

\section{Record mechanisms}\label{sec:mech}

The census records decompose into a small move inventory [MV for the instances;
observational as a classification]: linear factors $x - a$ for available constants
$a$; block translation at one gate per unit shift; multiplication by the input $x$
(one gate, adjoins $0$); and difference-of-squares splittings $B^2 - A^2 = (B-A)(B+A)$
at two gates given $A$ and $B$. The depth-$5$ record
$x^2 - (x^2 - 2)^2 = -(x{-}1)(x{+}1)(x{-}2)(x{+}2)$ is the fundamental instance: its
inner map $C(x) = x^2 - 2$ is the Chebyshev-conjugate doubling map
($h(z) = z + 1/z$ satisfies $h(z)^2 - 2 = h(z^2)$), i.e.\ the record is the integer
shadow of exactly the iteration that makes the \emph{real} analogue of the conjecture
false. Over $\Z$, that factory is sterile beyond one splitting, which the next
section proves.

\section{Stall theorems}\label{sec:stall}

\begin{lemma}[escape]\label{lem:escape}
Let $h = x^d + \sum_{i<d} a_i x^i \in \Z[x]$, $d \ge 2$, and
$R = 1 + \sum_{i<d} |a_i|$. If $|x| \ge R + 1$ then $|h(x)| \ge |x| + 1$. \emph{[D]}
\end{lemma}
\begin{proof}
$\sum_{i<d} |a_i|\,|x|^i \le (R-1)|x|^{d-1}$ for $|x| \ge 1$, so
$|h(x)| \ge |x|^{d-1}(|x| - (R-1)) \ge 2|x|^{d-1} \ge 2|x| \ge |x| + 1$ for
$|x| \ge R + 1$.
\end{proof}

\begin{theorem}[monic stall; D]\label{thm:stall}
Fix a monic $h \in \Z[x]$ of degree $\ge 2$. There is a constant $Z(h)$ such that
for every $k \ge 1$, both $h^{\circ k}(x) - x$ and the difference-of-squares towers
$G_k = h^{\circ (k-1)}(x)^2 - h^{\circ k}(x)^2$ have at most $Z(h)$ distinct integer
roots. One may take $Z(h) = \#([-M, M] \cap \Z)$ where
$M = \max(R+1, \max |a|)$ over the integer solutions $a$ of $h(y) = \pm y$.
\end{theorem}
\begin{proof}
Integer roots of $h^{\circ k}(x) - x$ are periodic points of $h$; by
Lemma~\ref{lem:escape} all integer periodic points lie in $[-R-1, R+1]$. For $G_k$,
factor $G_k = (h^{\circ(k-1)} - h^{\circ k})(h^{\circ(k-1)} + h^{\circ k})$: an
integer root $x$ has $y = h^{\circ(k-1)}(x)$ solving $h(y) = \pm y$, so $x$ lies in
the increasing union $\bigcup_j (h^{\circ j})^{-1}(A) \cap \Z$ for the finite set $A$
of such $y$. If $|x| > M$ then $|h^{\circ j}(x)| \ge |x| + j > \max_{a \in A} |a|$
for all $j$ by Lemma~\ref{lem:escape} and induction, so no element of the union
exceeds $M$ in absolute value; the union is contained in a fixed finite set and
stabilizes.
\end{proof}

\begin{proposition}[the Chebyshev core; D, MV spot-checks]\label{prop:cheb}
For $C(x) = x^2 - 2$: the integer periodic points are exactly $\{2, -1\}$; for every
$k \ge 1$ the polynomial $C^{\circ k}(x) - x$ has exactly $2$ distinct integer roots
while having $2^k$ distinct real roots and $\tau \le 2k + 2$; and the towers $G_k$
have root set $\{\pm 1, \pm 2\}$ for $k = 1$ and exactly $\{0, \pm 1, \pm 2\}$ for
every $k \ge 2$.
\end{proposition}
\begin{proof}
Orbits: $0 \mapsto -2 \mapsto 2 \mapsto 2$ and $1 \mapsto -1 \mapsto -1$; by
Lemma~\ref{lem:escape} (with $R = 3$; direct check for $|x| \ge 3$) all other integer
orbits escape, so the only cycles are the fixed points $2$ and $-1$, giving the first
two claims (the real-root count is the classical Chebyshev conjugacy). For $G_k$:
$h(y) = y$ has integer solutions $\{2, -1\}$ and $h(y) = -y$ has $\{-2, 1\}$;
integer preimages are $C^{-1}(2) = \{\pm 2\}$, $C^{-1}(-1) = \{\pm 1\}$,
$C^{-1}(-2) = \{0\}$, $C^{-1}(1) = C^{-1}(0) = \emptyset$, so the preimage closure
of $\{\pm 1, \pm 2\}$ is $\{0, \pm 1, \pm 2\}$ and is stable.
\end{proof}

Theorem~\ref{thm:stall} closes a construction class: no family built by iterating one
fixed monic map can witness superpolynomial $z$ versus $\tau$; a refutation of the
conjecture, if one exists, must build fresh constants or mix maps, paying gates for
them. That is precisely the regime the census measures.

\section{The quadratic family}\label{sec:family}

We measured all tower shapes ($h^{\circ k}(x) - x$ and
$h^{\circ i}(x)^2 - h^{\circ j}(x)^2$, $0 \le i < j \le 4$) across
$h_c = x^2 - c$ for $1 \le c \le 200$, in exact arithmetic. Root finding here cannot
use trailing-coefficient divisors (those constants reach $c^{2^k} \sim 10^{18}$);
instead, by Lemma~\ref{lem:escape}, every integer root of such a shape lies in
$[-(c+1), c+1]$, and direct evaluation over that window is exact. The two finders
agree on every shape where both are feasible ($c \le 10$, $j \le 2$) [MV].

\begin{theorem}[family census; MV]\label{thm:family}
For $1 \le c \le 200$ and the shapes above: the maximum number of distinct integer
roots is $5$, attained only at $c = 2$ (the set $\{0, \pm1, \pm2\}$ of
Proposition~\ref{prop:cheb}). Nonzero yields occur exactly on two series:
$c = m(m+1)$, yield $4$ for $m \ge 2$ with root set $\{\pm m, \pm(m+1)\}$ (from the
fixed points $\{m{+}1, -m\}$ and anti-fixed points $\{m, -m{-}1\}$), and
$c = m^2 + m + 1$, yield $4$ with the same root sets, produced by the genuine
$2$-cycles $m \mapsto -m-1 \mapsto m$ (for $c = 1$: yield $3$, root set
$\{0, \pm 1\}$, from the cycle $0 \leftrightarrow -1$).
\end{theorem}

The second series refuted our committed prediction that only the fixed-point series
yields roots. It is, however, the \emph{last} such surprise available to this
construction class:

\begin{proposition}[cycle ceiling; D, classical]\label{prop:cycles}
For any $f \in \Z[x]$, every cycle of $f$ in $\Z$ has length at most $2$.
Consequently fixed points, anti-fixed points and $2$-cycles are the complete
periodic inventory harvestable by tower shapes over any integer polynomial map.
\end{proposition}
\begin{proof}
For integers $a \ne b$, $(a - b) \mid (f(a) - f(b))$. Around a cycle
$a_0 \mapsto a_1 \mapsto \dots \mapsto a_r = a_0$ of distinct points, the
differences $d_i = a_{i+1} - a_i$ satisfy $d_i \mid d_{i+1}$ cyclically, so all
$|d_i|$ equal one value $d > 0$, and the cycle is a closed walk on the arithmetic
line $a_0 + d\,\Z$ with steps $\pm d$. If two consecutive steps have opposite signs
then $a_{i+2} = a_i$, which forces $r = 2$; if all steps have the same sign the walk
never closes. Hence $r \le 2$. The fact is classical in the polynomial-cycles
literature (see Narkiewicz \cite{Narkiewicz1995}); the proof is included to keep the
paper self-contained.
\end{proof}

Yield per gate across the family is maximized at $c = 2$ and decreases in $c$ (the
yield is bounded by $5$ while building $c$ costs $\tau(c) \sim \log c$ gates,
measured exactly against our anchored integer census): within this family, the
``parameterized loophole'' left open by Theorem~\ref{thm:stall} is empty [MV for
$c \le 200$; C beyond].

\section{Limits, and what would count as progress}\label{sec:limits}

Nothing above bears on the conjecture asymptotically: censuses to depth $7$ cannot
distinguish polynomial from superpolynomial growth, and our theorems close
construction classes rather than open the bounding problem. The honest state of the
measured landscape is: every mechanism observed or excluded so far pays $\Theta(1)$
gates per root, the growth function's first structure (plateaus at $4$ and $7$) is
the visible cost of constant-building, and the geometric-progression family remains
the best known factory, linear-rate, exactly as when Rojas noted that the conjecture
fails for $\kappa < 1$ \cite{Rojas2003}.

Concrete open items this record makes precise:
(i) the exact value of $\zmax(8)$ (now known only to be $\ge 6$), including whether
an $8$-gate seven-rooter exists, and the seven-root threshold: finite, SAT-shaped
questions in the style of exact SLP synthesis over GF(2)
\cite{FuhsSchneiderKamp2010};
(ii) extend the integer frontier past Markstr\"om's length $11$ \cite{Markstrom2014},
including his monotonicity question for $\tau(n!)$;
(iii) measure the growth of the valuation-spectrum bound
$s \le N_p(s) \le s(s+1)/2$ of \cite{Rojas2003}, where the true rate is stated as
open;
(iv) the census's $2$-adic instrumentation records that all depth-$5$--$7$ record
polynomials concentrate their roots on the valuations $\{0, 1\}$ [MV], consistent
with the pressure sitting on the digit-conjecture side of Rojas' reduction [C].

\section*{Data and code availability}

All code, hypotheses, run artifacts, verdicts and the wiki are committed at
\texttt{github.com/fsantibanezleal/CAOS\_RESEARCH} under
\texttt{problems/computation-complexity/tau-conjecture/} (experiments
\texttt{EXP-001}--\texttt{EXP-005}; package \texttt{tclib} with its test suite). All
computations are exact integer arithmetic in Python 3.13 standard library; total
compute for every result in this paper is under five hours on one desktop CPU
(24 cores available, single-threaded runs; peak memory $\sim$6\,GB).

\begin{thebibliography}{99}\small

\bibitem{ShubSmale1995} M.~Shub, S.~Smale,
\emph{On the intractability of Hilbert's Nullstellensatz and an algebraic version of
``NP $\ne$ P?''}, Duke Math.\ J.\ 81(1):47--54, 1995.
\href{https://doi.org/10.1215/S0012-7094-95-08105-8}{doi:10.1215/S0012-7094-95-08105-8}.

\bibitem{Smale1998} S.~Smale, \emph{Mathematical problems for the next century},
Math.\ Intelligencer 20(2):7--15, 1998.
\href{https://doi.org/10.1007/BF03025291}{doi:10.1007/BF03025291}.

\bibitem{BCSS1998} L.~Blum, F.~Cucker, M.~Shub, S.~Smale, \emph{Complexity and Real
Computation}, Springer, 1998.

\bibitem{Koiran2004} P.~Koiran, \emph{Valiant's model and the cost of computing
integers}, Comput.\ Complexity 13(3--4):131--146, 2004.
\href{https://doi.org/10.1007/s00037-004-0186-2}{doi:10.1007/s00037-004-0186-2}.

\bibitem{Burgisser2009} P.~B\"urgisser, \emph{On defining integers and proving
arithmetic circuit lower bounds}, Comput.\ Complexity 18(1):81--103, 2009.
\href{https://doi.org/10.1007/s00037-009-0260-x}{doi:10.1007/s00037-009-0260-x}.

\bibitem{Burgisser2024} P.~B\"urgisser, \emph{Completeness classes in algebraic
complexity theory}, arXiv:2406.06217, 2024.

\bibitem{Markstrom2014} K.~Markstr\"om, \emph{The straight line complexity of small
factorials and primorials}, INTEGERS 14, \#A67, 2014. arXiv:1306.3091.

\bibitem{Rojas2003} J.~M.~Rojas, \emph{A direct ultrametric approach to additive
complexity and the Shub--Smale tau conjecture}, arXiv:math/0304100, 2003.

\bibitem{Cheng2003} Q.~Cheng, \emph{Straight-line programs and torsion points on
elliptic curves}, Comput.\ Complexity 12(3--4):150--161, 2003.
\href{https://doi.org/10.1007/s00037-003-0180-0}{doi:10.1007/s00037-003-0180-0}.

\bibitem{Cheng2004} Q.~Cheng, \emph{On the ultimate complexity of factorials},
Theoret.\ Comput.\ Sci.\ 326(1--3):419--429, 2004.
\href{https://doi.org/10.1016/j.tcs.2004.07.032}{doi:10.1016/j.tcs.2004.07.032}.

\bibitem{Shamir1979} A.~Shamir, \emph{Factoring numbers in $O(\log n)$ arithmetic
steps}, Inform.\ Process.\ Lett.\ 8(1):28--31, 1979.
\href{https://doi.org/10.1016/0020-0190(79)90087-5}{doi:10.1016/0020-0190(79)90087-5}.

\bibitem{Lipton1994} R.~J.~Lipton, \emph{Straight-line complexity and integer
factorization}, ANTS-I, LNCS 877, 1994.
\href{https://doi.org/10.1007/3-540-58691-1_50}{doi:10.1007/3-540-58691-1\_50}.

\bibitem{Koiran2011} P.~Koiran, \emph{Shallow circuits with high-powered inputs},
Innovations in Computer Science, 2011. arXiv:1004.4960.

\bibitem{KPTT2015} P.~Koiran, N.~Portier, S.~Tavenas, S.~Thomass\'e, \emph{A
tau-conjecture for Newton polygons}, Found.\ Comput.\ Math.\ 15:185--197, 2015.
arXiv:1308.2286.

\bibitem{Hrubes2013} P.~Hrube\v{s}, \emph{On the real $\tau$-conjecture and the
distribution of complex roots}, Theory of Computing 9(10):403--411, 2013.
\href{https://doi.org/10.4086/toc.2013.v009a010}{doi:10.4086/toc.2013.v009a010}.

\bibitem{BriquelBurgisser2020} I.~Briquel, P.~B\"urgisser, \emph{The real
tau-conjecture is true on average}, Random Struct.\ Algorithms 57(2):279--303, 2020.
\href{https://doi.org/10.1002/rsa.20926}{doi:10.1002/rsa.20926}.

\bibitem{Dutta2021} P.~Dutta, \emph{Real $\tau$-conjecture for sum-of-squares: a
unified approach to lower bound and derandomization}, CSR 2021.
\href{https://doi.org/10.1007/978-3-030-79416-3_5}{doi:10.1007/978-3-030-79416-3\_5}.

\bibitem{BBDM2026} S.~Bhattacharjee, M.~Bl\"aser, P.~Dutta, S.~Mukherjee,
\emph{Exponential lower bounds via exponential sums}, ICALP 2024; extended version
arXiv:2601.00387, 2026.

\bibitem{Burgisser2026} P.~B\"urgisser, \emph{Intractability of Hilbert's
Nullstellensatz implies algebraic hardness of the permanent}, arXiv:2606.25121, 2026.

\bibitem{FuhsSchneiderKamp2010} C.~Fuhs, P.~Schneider-Kamp, \emph{Synthesizing
shortest linear straight-line programs over GF(2) using SAT}, SAT 2010, LNCS 6175.
\href{https://doi.org/10.1007/978-3-642-14186-7_8}{doi:10.1007/978-3-642-14186-7\_8}.

\bibitem{Narkiewicz1995} W.~Narkiewicz, \emph{Polynomial Mappings}, Lecture Notes in
Mathematics 1600, Springer, 1995.

\end{thebibliography}

\end{document}
